

% This is a simple sample document.  For more complicated documents take a look in the exercise tab. Note that everything that comes after a % symbol is treated as comment and ignored when the code is compiled.

\documentclass[fleqn]{report} % \documentclass{} is the first command in any LaTeX code.  It is used to define what kind of document you are creating such as an article or a book, and begins the document preamble

\usepackage{amsmath} % \usepackage is a command that allows you to add functionality to your LaTeX code
\usepackage{amsfonts}
\usepackage{algorithm}
\usepackage{algpseudocode}
\usepackage{graphicx}
\usepackage{xcolor}
\newcommand\ddfrac[2]{\frac{\displaystyle #1}{\displaystyle #2}}
%\usepackage[linesnumbered,ruled,vlined]{algorithm2e}



\title{Technical report\\Sparse Principal Component Solver} % Sets article title
\author{Dr. Mamadou THIAO} % Sets authors name
\date{\today} % Sets date for date compiled
\newtheorem{thm}{Proposition}[subsection]
\newenvironment{proof}{\paragraph{Proof:}}{\hfill$\square$}
\newtheorem{crly}{Corollary}[subsection]

% The preamble ends with the command \begin{document}
\begin{document} % All begin commands must be paired with an end command somewhere
    \maketitle % creates title using information in preamble (title, author, date)
    \tableofcontents
    \section{Sparse Principal Component Problem} % creates a section
A formulation of the Sparse Principal Component Problem is as follows
\begin{equation}
		\label{Constrained_spca}
		(\text{SPCA}_{\Sigma}^{k})\;
		\left\{
		\begin{array}{ll}\max & x^{T}\Sigma x\\
			s.t.&x^{T}x = 1,\\
			&\left\|x\right\|_0 \leq k,\\
			&\;x\in\mathbb{R}^n,\\
		\end{array} 
		\right.
    \end{equation}
		where the integer parameter $1 \leq k\leq n $ controls the sparsity of the solution.
    \section{Extracting several sparse components} % creates a section
A well know technique to extract more than one sparse component is by deflating the covariance matrix at each round. see Mackey et al.
\begin{algorithm}[H]
        \caption{Generalized Deflation Method for Sparse PCA}
        \label{alg:gdmspca}
        \begin{algorithmic}[1] % [1] adds line numbering
	  \State Given: covariance matrix  $\Sigma_0$, number of rounds $r$, $\left\{k_1,...,k_r\right\}\subset\left\{1, ..., n\right\}$.
	\State $B_0 := I$,
            \For{$t=1, ..., r$}
		\State $x_t := \underset{x^{T}B_{t-1}x = 1, \left\|x\right\|_0 \leq k}{\text{argmax}}\; x^{T}\Sigma_{t-1} x$
		\State  $q_t :=B_{t-1}x_t$
		\State $\Sigma_{t} := (I-q_t^{T}q_t)\Sigma_{t-1}(I-q_t^{T}q_t)$
		\State $B_{t} := B_{t-1}(I-q_t^{T}q_t)$
		\State $x_{t} := x_{t} \slash\left\| x_{t}\right\|_2$
	 \EndFor
            \State \Return $\left\{x_1,...,x_r\right\}$
        \end{algorithmic}
    \end{algorithm}
An approach for solving point 4: in Algorithm~\ref{alg:gdmspca} is to first replace $x^{T}B_{t-1}x = 1$ by $x^{T}(B_{t-1} + \epsilon I)x = 1$ and use Woodbury matrix identity to transform it in the form of \ref{Constrained_spca}. To be detailed!
\section{Implemented methods} 
\begin{itemize}
\item Dca method from Thiao et al.
\item Gspca forward and backward methods from Moghaddam et al.
\item We will implement all the sparse principal component methods and deflation techniques and incorporate them into the solver.
\end{itemize}
\begin{thebibliography}{99} % The '9' handles up to 9 references

\bibitem{Mackey2010}
Thiao, M., Pham Dinh, T., \& Le Thi, H. A. L.: \textit{A DC Programming Approach for Sparse Eigenvalue Problem}. Proceedings of the 27th International Conference on Machine Learning (2010).

\bibitem{Mackey2008}
Mackey, L.: \textit{Deflation methods for sparse PCA.} Advances in Neural Information Processing Systems (2008), 21:1017–1024.

\bibitem{Moghaddam2006}
Moghaddam, B., Weiss, Y., \& Avidan, S.: \textit{Spectral bounds for sparse PCA: Exact and greedy
algorithms.} Advances in Neural Information Processing Systems (2006), 18.

\end{thebibliography}
\end{document} % This is the end of the document
